\documentclass[10pt,a4paper]{amsart}
\usepackage{amsmath,amssymb,amsthm,graphicx}
% Retained minimal preamble for the Stacks Project chapter more-morphisms.tex (More on Morphisms).
% Verbatim excerpt of the chapter preamble: Xy-pic setup (native preamble lines 14-18) and the native
% theorem declarations with the mathematical macros (native preamble lines 154-217). Copied mechanically;
% nothing added, removed or redefined. The wrapper class and the amsmath/amssymb/amsthm/graphicx base
% packages are supplied by the assembler, so the class/externaldocument part of the native preamble is omitted.
% For commutative diagrams we use Xy-pic
\usepackage[all]{xy}

% We use 2cell for 2-commutative diagrams.
\xyoption{2cell}
% Theorem environments.
%
\theoremstyle{plain}
\newtheorem{theorem}[subsection]{Theorem}
\newtheorem{proposition}[subsection]{Proposition}
\newtheorem{lemma}[subsection]{Lemma}

\theoremstyle{definition}
\newtheorem{definition}[subsection]{Definition}
\newtheorem{example}[subsection]{Example}
\newtheorem{exercise}[subsection]{Exercise}
\newtheorem{situation}[subsection]{Situation}

\theoremstyle{remark}
\newtheorem{remark}[subsection]{Remark}
\newtheorem{remarks}[subsection]{Remarks}

\numberwithin{equation}{subsection}

% Macros
%
\def\lim{\mathop{\mathrm{lim}}\nolimits}
\def\colim{\mathop{\mathrm{colim}}\nolimits}
\def\Spec{\mathop{\mathrm{Spec}}}
\def\Hom{\mathop{\mathrm{Hom}}\nolimits}
\def\Ext{\mathop{\mathrm{Ext}}\nolimits}
\def\SheafHom{\mathop{\mathcal{H}\!\mathit{om}}\nolimits}
\def\SheafExt{\mathop{\mathcal{E}\!\mathit{xt}}\nolimits}
\def\Sch{\mathit{Sch}}
\def\Mor{\mathop{\mathrm{Mor}}\nolimits}
\def\Ob{\mathop{\mathrm{Ob}}\nolimits}
\def\Sh{\mathop{\mathit{Sh}}\nolimits}
\def\NL{\mathop{N\!L}\nolimits}
\def\CH{\mathop{\mathrm{CH}}\nolimits}
\def\proetale{{pro\text{-}\acute{e}tale}}
\def\etale{{\acute{e}tale}}
\def\QCoh{\mathit{QCoh}}
\def\Ker{\mathop{\mathrm{Ker}}}
\def\Im{\mathop{\mathrm{Im}}}
\def\Coker{\mathop{\mathrm{Coker}}}
\def\Coim{\mathop{\mathrm{Coim}}}

% Boxtimes
%
\DeclareMathSymbol{\boxtimes}{\mathbin}{AMSa}{"02}

%
% Macros for moduli stacks/spaces
%
\def\QCohstack{\mathcal{QC}\!\mathit{oh}}
\def\Cohstack{\mathcal{C}\!\mathit{oh}}
\def\Spacesstack{\mathcal{S}\!\mathit{paces}}
\def\Quotfunctor{\mathrm{Quot}}
\def\Hilbfunctor{\mathrm{Hilb}}
\def\Curvesstack{\mathcal{C}\!\mathit{urves}}
\def\Polarizedstack{\mathcal{P}\!\mathit{olarized}}
\def\Complexesstack{\mathcal{C}\!\mathit{omplexes}}
% \Pic is the operator that assigns to X its picard group, usage \Pic(X)
% \Picardstack_{X/B} denotes the Picard stack of X over B
% \Picardfunctor_{X/B} denotes the Picard functor of X over B
\def\Pic{\mathop{\mathrm{Pic}}\nolimits}
\def\Picardstack{\mathcal{P}\!\mathit{ic}}
\def\Picardfunctor{\mathrm{Pic}}
\def\Deformationcategory{\mathcal{D}\!\mathit{ef}}
\usepackage{cleveref}
\crefname{originalreference}{original reference}{original references}
\Crefname{originalreference}{Original reference}{Original references}
\makeatletter
\newlabel{cohomology-lemma-projection-formula-closed-immersion}{{0B55}{source}{}{sourcepointer-1}{}}
\newlabel{cohomology-lemma-projection-formula-closed-immersion@cref}{{[originalreference][0][]0B55}{[1][0][]source}{}{}{}}
\makeatother
\crefname{originalreference}{original reference}{original references}
\Crefname{originalreference}{Original reference}{Original references}
\makeatletter
\newlabel{more-algebra-lemma-check-pseudo-coherent-modulo-nilpotent}{{0H76}{source}{}{sourcepointer-2}{}}
\newlabel{more-algebra-lemma-check-pseudo-coherent-modulo-nilpotent@cref}{{[originalreference][0][]0H76}{[1][0][]source}{}{}{}}
\makeatother
\crefname{originalreference}{original reference}{original references}
\Crefname{originalreference}{Original reference}{Original references}
\makeatletter
\newlabel{perfect-lemma-affine-compare-bounded}{{06Z0}{source}{}{sourcepointer-3}{}}
\newlabel{perfect-lemma-affine-compare-bounded@cref}{{[originalreference][0][]06Z0}{[1][0][]source}{}{}{}}
\makeatother
\crefname{originalreference}{original reference}{original references}
\Crefname{originalreference}{Original reference}{Original references}
\makeatletter
\newlabel{perfect-lemma-pseudo-coherent}{{08E5}{source}{}{sourcepointer-4}{}}
\newlabel{perfect-lemma-pseudo-coherent@cref}{{[originalreference][0][]08E5}{[1][0][]source}{}{}{}}
\makeatother
\crefname{originalreference}{original reference}{original references}
\Crefname{originalreference}{Original reference}{Original references}
\makeatletter
\newlabel{perfect-lemma-pseudo-coherent-affine}{{08E7}{source}{}{sourcepointer-5}{}}
\newlabel{perfect-lemma-pseudo-coherent-affine@cref}{{[originalreference][0][]08E7}{[1][0][]source}{}{}{}}
\makeatother
\crefname{originalreference}{original reference}{original references}
\Crefname{originalreference}{Original reference}{Original references}
\makeatletter
\newlabel{perfect-lemma-quasi-coherence-direct-image}{{08D5}{source}{}{sourcepointer-6}{}}
\newlabel{perfect-lemma-quasi-coherence-direct-image@cref}{{[originalreference][0][]08D5}{[1][0][]source}{}{}{}}
\makeatother
\crefname{originalreference}{original reference}{original references}
\Crefname{originalreference}{Original reference}{Original references}
\makeatletter
\newlabel{perfect-lemma-quasi-coherence-tensor-product}{{08DX}{source}{}{sourcepointer-7}{}}
\newlabel{perfect-lemma-quasi-coherence-tensor-product@cref}{{[originalreference][0][]08DX}{[1][0][]source}{}{}{}}
\makeatother
\crefname{originalsection}{original section}{original sections}
\Crefname{originalsection}{Original section}{Original sections}
\makeatletter
\newlabel{section-thickenings}{{04EW}{source}{}{sourcepointer-8}{}}
\newlabel{section-thickenings@cref}{{[originalsection][0][]04EW}{[1][0][]source}{}{}{}}
\makeatother
\title{Pseudo-coherence ascends along a finite order thickening of schemes}
\author{The Stacks Project Authors}
\date{stacks-project:more-morphisms.tex:local-snapshot-20261011; GNU FDL 1.2 or later}
\begin{document}
\maketitle
\noindent\textit{Editorial note.} Counted span is the author's own proof environment, source lines 20743-20793, immediately following the statement at 20736-20741.\par\medskip
\noindent\textit{Editorial note.} Seven cross-chapter pointers (Cohomology, More on Algebra, Derived Categories of Schemes) plus the chapter's own thickenings section are recorded in assembly.external\_references.\par\medskip
\noindent\textit{Editorial note.} Same-name pair note: the spaces side of this label delegates its core case to this proof, so the pair is counted once, here.\par\medskip
\noindent\textit{Allowed standard prerequisites (editorial).}\begin{itemize}
\item Section on thickenings. chapter Section section-thickenings, cited at proof line 20746. Reduces the general finite order thickening to a first order one.
\item Pseudo-coherence modulo a nilpotent. More on Algebra, lemma more-algebra-lemma-check-pseudo-coherent-modulo-nilpotent, cited in the proof. Is the algebraic criterion that lets pseudo-coherence be tested after reduction modulo the nilpotent ideal.
\item Affine comparison for bounded complexes. Derived Categories of Schemes, lemma perfect-lemma-affine-compare-bounded, cited at proof lines 20784-20785. Represents the object by a complex of modules over A' so that the ring-theoretic criterion applies.
\item Quasi-coherence of direct images and tensor products. Derived Categories of Schemes, lemmas perfect-lemma-quasi-coherence-direct-image and perfect-lemma-quasi-coherence-tensor-product, cited in the proof. Establishes that K' has quasi-coherent cohomology sheaves, the first half of the proof.
\item Projection formula for a closed immersion. Cohomology, lemma cohomology-lemma-projection-formula-closed-immersion, cited in the proof. Computes the restriction of K' to X in the short exact sequence of the thickening.
\end{itemize}\medskip
\ifdefined\hypertarget\hypertarget{sourcepointer-1}{}\fi
\noindent\textit{Source reference.} Original reference 0B55: focused-seeds/geometry/S9-morph-thickening-pseudo-coherent-schemes/ref-tags.txt, source lines 4--4. Navigation pointer only: the counted proof cites this statement by \textbackslash{}ref as a standing input; the deduction the counted proof has to make is carried out inside the counted span, and no part of that deduction is delegated to this pointer.\par\medskip
\ifdefined\hypertarget\hypertarget{sourcepointer-2}{}\fi
\noindent\textit{Source reference.} Original reference 0H76: focused-seeds/geometry/S9-morph-thickening-pseudo-coherent-schemes/ref-tags.txt, source lines 5--5. Navigation pointer only: the counted proof cites this statement by \textbackslash{}ref as a standing input; the deduction the counted proof has to make is carried out inside the counted span, and no part of that deduction is delegated to this pointer.\par\medskip
\ifdefined\hypertarget\hypertarget{sourcepointer-3}{}\fi
\noindent\textit{Source reference.} Original reference 06Z0: focused-seeds/geometry/S9-morph-thickening-pseudo-coherent-schemes/ref-tags.txt, source lines 6--6. Navigation pointer only: the counted proof cites this statement by \textbackslash{}ref as a standing input; the deduction the counted proof has to make is carried out inside the counted span, and no part of that deduction is delegated to this pointer.\par\medskip
\ifdefined\hypertarget\hypertarget{sourcepointer-4}{}\fi
\noindent\textit{Source reference.} Original reference 08E5: focused-seeds/geometry/S9-morph-thickening-pseudo-coherent-schemes/ref-tags.txt, source lines 7--7. Navigation pointer only: the counted proof cites this statement by \textbackslash{}ref as a standing input; the deduction the counted proof has to make is carried out inside the counted span, and no part of that deduction is delegated to this pointer.\par\medskip
\ifdefined\hypertarget\hypertarget{sourcepointer-5}{}\fi
\noindent\textit{Source reference.} Original reference 08E7: focused-seeds/geometry/S9-morph-thickening-pseudo-coherent-schemes/ref-tags.txt, source lines 8--8. Navigation pointer only: the counted proof cites this statement by \textbackslash{}ref as a standing input; the deduction the counted proof has to make is carried out inside the counted span, and no part of that deduction is delegated to this pointer.\par\medskip
\ifdefined\hypertarget\hypertarget{sourcepointer-6}{}\fi
\noindent\textit{Source reference.} Original reference 08D5: focused-seeds/geometry/S9-morph-thickening-pseudo-coherent-schemes/ref-tags.txt, source lines 9--9. Navigation pointer only: the counted proof cites this statement by \textbackslash{}ref as a standing input; the deduction the counted proof has to make is carried out inside the counted span, and no part of that deduction is delegated to this pointer.\par\medskip
\ifdefined\hypertarget\hypertarget{sourcepointer-7}{}\fi
\noindent\textit{Source reference.} Original reference 08DX: focused-seeds/geometry/S9-morph-thickening-pseudo-coherent-schemes/ref-tags.txt, source lines 10--10. Navigation pointer only: the counted proof cites this statement by \textbackslash{}ref as a standing input; the deduction the counted proof has to make is carried out inside the counted span, and no part of that deduction is delegated to this pointer.\par\medskip
\ifdefined\hypertarget\hypertarget{sourcepointer-8}{}\fi
\noindent\textit{Source reference.} Original section 04EW: sources/stacks-more-morphisms/pinned-original-source.tex, source lines 30--30. Navigation pointer only: the counted proof cites this statement by \textbackslash{}ref as a standing input; the deduction the counted proof has to make is carried out inside the counted span, and no part of that deduction is delegated to this pointer.\par\medskip
\begin{lemma}
\label{lemma-thickening-pseudo-coherent}
Let $i : X \to X'$ be a finite order thickening of schemes. Let
$K' \in D(\mathcal{O}_{X'})$ be an object such that
$K = Li^*K'$ is pseudo-coherent. Then $K'$ is pseudo-coherent.
\end{lemma}

\begin{proof}
We first prove $K'$ has quasi-coherent cohomology sheaves.
To do this, we may reduce to the case of a first order thickening, see
Section \ref{section-thickenings}. Let $\mathcal{I} \subset \mathcal{O}_{X'}$
be the quasi-coherent sheaf of ideals cutting out $X$.
Tensoring the short exact sequence
$$
0 \to \mathcal{I} \to \mathcal{O}_{X'} \to i_*\mathcal{O}_X \to 0
$$
with $K'$ we obtain a distinguished triangle
$$
K' \otimes_{\mathcal{O}_{X'}}^\mathbf{L} \mathcal{I}
\to K' \to
K' \otimes_{\mathcal{O}_{X'}}^\mathbf{L} i_*\mathcal{O}_X
\to
(K' \otimes_{\mathcal{O}_{X'}}^\mathbf{L} \mathcal{I})[1]
$$
Since $i_* = Ri_*$ and since we may view $\mathcal{I}$
as a quasi-coherent $\mathcal{O}_X$-module (as we have a first
order thickening) we may rewrite this as
$$
i_*(K \otimes_{\mathcal{O}_X}^\mathbf{L} \mathcal{I})
\to K' \to
i_*K \to
i_*(K \otimes_{\mathcal{O}_X}^\mathbf{L} \mathcal{I})[1]
$$
Please use Cohomology, Lemma
\ref{cohomology-lemma-projection-formula-closed-immersion}
to identify the terms. Since $K$ is in
$D_\QCoh(\mathcal{O}_X)$ we conclude that
$K'$ is in $D_\QCoh(\mathcal{O}_{X'})$; this uses
Derived Categories of Schemes, Lemmas
\ref{perfect-lemma-pseudo-coherent},
\ref{perfect-lemma-quasi-coherence-tensor-product}, and
\ref{perfect-lemma-quasi-coherence-direct-image}.

\medskip\noindent
Assume $K'$ is in $D_\QCoh(\mathcal{O}_{X'})$.
The question is local on $X'$ hence we may assume $X'$ is affine.
Say $X' = \Spec(A')$ and $X = \Spec(A)$ with $A = A'/I$ and $I$ nilpotent.
Then $K'$ comes from an object $M' \in D(A')$, see
Derived Categories of Schemes, Lemma
\ref{perfect-lemma-affine-compare-bounded}.
Thus $M = M' \otimes_{A'}^\mathbf{L} A$ is a pseudo-coherent
object of $D(A)$ by Derived Categories of Schemes,
Lemma \ref{perfect-lemma-pseudo-coherent-affine} and our assumption on $K$.
It follows that $M'$ is pseudo-coherent by
More on Algebra, Lemma
\ref{more-algebra-lemma-check-pseudo-coherent-modulo-nilpotent}.
This finishes the proof.
\end{proof}

\end{document}
